\section*{Chapter \ref{ch:b3:ocean-dynamics} --- Ocean Dynamics}

\begin{solution}{exo:b3:ocean-dynamics:1}
$f = \qty{9.37e-5}{s^{-1}}$; $\tau/(\rho f) = 0.12/(1025\times\num{9.37e-5})
= \qty{1.25}{m^2/s}$, to the right of the wind: south. Across
\qty{1000}{km}, \qty{1.25e6}{m^3/s}, \qty{1.25}{\sverdrup}.
\end{solution}

\begin{solution}{exo:b3:ocean-dynamics:2}
$f = \qty{1.031e-4}{s^{-1}}$; $D_E = \pi\sqrt{2\times0.05/\num{1.031e-4}}
= \qty{98}{m}$. The current there is $\mathrm e^{-\pi} = \num{0.043}$ of the
surface current, flowing the opposite way.
\end{solution}

\begin{solution}{exo:b3:ocean-dynamics:3}
$\beta = 2\Omega\cos\ang{30}/a = \qty{1.98e-11}{m^{-1}.s^{-1}}$; $V =
\num{-1e-7}/(1025\times\num{1.98e-11}) = \qty{-4.92}{m^2/s}$; across
\qty{5000}{km}, \qty{-24.6}{\sverdrup}: southward.
\end{solution}

\begin{solution}{exo:b3:ocean-dynamics:4}
The two constituents come back into phase after $1/(1/12 - 1/12.4206) =
\qty{354}{h}$, \qty{14.77}{days}, half a synodic month.
\end{solution}

\begin{solution}{exo:b3:ocean-dynamics:5}
At \ang{40} N, $f = \qty{9.37e-5}{s^{-1}}$, $\beta =
\qty{1.75e-11}{m^{-1}.s^{-1}}$, $\tau_x = \qty{0.05}{Pa}$ and
$\partial\tau_x/\partial y = 0.1/\num{1.11e6} = \qty{9.0e-8}{Pa/m}$. The
first term, $-\num{9.0e-8}/(1025\times\num{9.37e-5}) =
\qty{-9.38e-7}{m/s}$; the second, $0.05\times\num{1.75e-11}/(1025\times
(\num{9.37e-5})^2) = \qty{+9.7e-8}{m/s}$. $w_E = \qty{-8.4e-7}{m/s}$,
\qty{-26.5}{m/yr}: the water is pushed down.
\end{solution}

\begin{solution}{exo:b3:ocean-dynamics:6}
$\delta_S = \num{2e-6}/\num{2e-11} = \qty{100}{km}$; speed $=
\num{30e6}/(\num{1e5}\times800) = \qty{0.375}{m/s}$.
\end{solution}

\begin{solution}{exo:b3:ocean-dynamics:7}
$\beta = \qty{1.98e-11}{m^{-1}.s^{-1}}$; $\delta_M = (10^4/\num{1.98e-11})^{1/3}
= \qty{79.6}{km}$, and the jet is $3.628\,\delta_M = \qty{289}{km}$ wide.
With $A_H$ eight times smaller both are halved, since $8^{1/3} = 2$.
\end{solution}

\begin{solution}{exo:b3:ocean-dynamics:8}
$f = \qty{8.57e-5}{s^{-1}}$; $v = (21.3 - 14.5)/(\num{8.57e-5}\times\num{8e4})
= \qty{0.99}{m/s}$, positive with $B$ on the right: looking from $A$ to $B$,
east, the current goes north. The sea surface stands
$6.8/9.81 = \qty{0.69}{m}$ higher at $B$; the water column at $B$ is the
lighter, warmer one, on the right of the current.
\end{solution}

\begin{solution}{exo:b3:ocean-dynamics:9}
Spring range $1.5K(1 + 0.46) = \qty{0.78}{m}$; neap range $1.5K(1 - 0.46)
= \qty{0.29}{m}$; ratio $1.46/0.54 = 2.7$.
\end{solution}

\begin{solution}{exo:b3:ocean-dynamics:10}
As in \cref{thm:b3:ocean-dynamics:kelvin}, $u = \pm(g/c)\eta$ for $F(x \mp
ct)$, and $fu = -g\,\partial\eta/\partial y$ becomes $\beta_0y(\pm g/c)G =
-gG'$, so $G = \exp(\mp\beta_0y^2/2c)$. Only the upper sign, the eastward
wave, is bounded on both sides of the equator. $\beta_0 =
\qty{2.29e-11}{m^{-1}.s^{-1}}$, $\sqrt{2.5/\num{2.29e-11}} = \qty{330}{km}$.
\end{solution}

\begin{solution}{exo:b3:ocean-dynamics:11}
(a)~Integrating $\beta v = f\,\partial w/\partial z$ from the bottom, $w = 0$,
to the top of the abyssal layer, $w = W_0$: $\beta V = fW_0$, and $f/\beta =
a\tan\varphi$. (b)~$V = \num{6.371e6}\times0.577\times10^{-7} =
\qty{0.368}{m^2/s}$, \qty{1.84}{\sverdrup} across \qty{5000}{km}, poleward.
(c)~From a deep western boundary current fed by the high-latitude sources;
the interior flows towards the sources, the boundary current away from
them.
\end{solution}

\begin{solution}{exo:b3:ocean-dynamics:12}
(a)~With $z$ upward from the floor, $A_zW'' = \mathrm if(W - W_g)$, the
pressure gradient being that of the geostrophic flow above. $W - W_g =
C\mathrm e^{-(1+\mathrm i)z/\delta}$ is bounded upward, and $W(0) = 0$ gives
$C = -W_g$. (b)~The transport defect is $\int_0^\infty(W - W_g)\,\mathrm dz
= -W_g\delta/(1+\mathrm i) = -\tfrac12\delta(1 - \mathrm i)W_g$, that is
$-\tfrac12\delta(u_g + v_g,\ v_g - u_g)$. The geostrophic flow being
non-divergent, $w = -\nabla\cdot(\text{defect}) = \tfrac12\delta(\partial_x
v_g - \partial_y u_g) = \tfrac12\delta\zeta_g$. (c)~$\delta =
\sqrt{0.02/10^{-4}} = \qty{14.1}{m}$; $r = 10^{-4}\times14.1/8000 =
\qty{1.77e-7}{s^{-1}}$, a spin-down time of \qty{65}{days}.
\end{solution}

\begin{solution}{pb:b3:ocean-dynamics:1}
\textbf{1.} $\tau_x = \qty{-0.1}{Pa}$ (easterly), 0 and \qty{+0.1}{Pa}
(westerly).

\textbf{2.} $f = \num{4.99e-5}$, \num{8.37e-5} and
\qty{1.117e-4}{s^{-1}}.

\textbf{3.} $V_E = -\tau_x/(\rho f)$: \qty{+1.96}{m^2/s} (north) at
\ang{20} N and \qty{-0.87}{m^2/s} (south) at \ang{50} N. They converge.

\textbf{4.} $D_E = \pi\sqrt{2\times0.05/\num{8.37e-5}} = \qty{109}{m}$.

\textbf{5.} $\tau_x = 0$ at \ang{35} N, so $w_E = -(1/\rho f)\,
\partial\tau_x/\partial y = -\pi\tau_0/(\rho fL) = \qty{-1.10e-6}{m/s}$,
\qty{-34.7}{m/yr}, downward.

\textbf{6.} $500/34.7 = \qty{14.4}{years}$.

\textbf{7.} $L = a\times\pi/6 = \qty{3336}{km}$; $\beta =
\qty{1.875e-11}{m^{-1}.s^{-1}}$.

\textbf{8.} $\operatorname{curl}_z\vec\tau = -\pi\tau_0/L =
\qty{-9.42e-8}{Pa/m}$.

\textbf{9.} $V = \num{-9.42e-8}/(1025\times\num{1.875e-11}) =
\qty{-4.90}{m^2/s}$, southward.

\textbf{10.} $4.90\times\num{6e6} = \qty{2.94e7}{m^3/s}$,
\qty{29.4}{\sverdrup} southward.

\textbf{11.} $\qty{4.9}{mm/s}$; \qty{1000}{km} in \qty{2.0e8}{s},
\qty{6.5}{years}.

\textbf{12.} $fw_E/\beta = \num{8.37e-5}\times(\num{-1.10e-6})/
\num{1.875e-11} = \qty{-4.90}{m^2/s}$: the same, because the Ekman
transport vanishes where $\tau_x = 0$. A column squashed at the rate $w_E$
keeps its \omterm{def:b3:atmospheric-dynamics:pv}{potential vorticity} $f/h$ only by moving to smaller $f$, south,
at the rate $\beta V_g = fw_E$.

\textbf{13.} A boundary layer can only be one in which the correction to
the interior decays away from the wall; with $\beta > 0$ that is
$\mathrm e^{-x/\delta_S}$, which decays eastward, from the western wall
only. In vorticity terms, friction in a northward current against an
eastern wall would remove positive vorticity, the wrong sign.

\textbf{14.} $\delta_S = \num{e-6}/\num{1.875e-11} = \qty{53.3}{km}$.

\textbf{15.} $A = \qty{4.90}{m^2/s}$, $X/\delta_S = 112.5$:
$4.90\times111.5 = \qty{546}{m^2/s}$, \qty{0.55}{m/s} through
\qty{1000}{m}.

\textbf{16.} Where $(X/\delta_S)\mathrm e^{-x/\delta_S} = 1$: $x =
\delta_S\ln(X/\delta_S) = \qty{252}{km}$.

\textbf{17.} $A(X - \delta_S - x_0) = 4.90\times(6000 - 53 - 252)\times
10^3 = \qty{27.9}{\sverdrup}$. The interior flow is still going south in
the strip next to the wall; the shortfall, $A(x_0 + \delta_S) =
\qty{1.5}{\sverdrup}$, is the interior transport that would have crossed
that strip.

\textbf{18.} $\delta_M = (\num{5e3}/\num{1.875e-11})^{1/3} = \qty{64.4}{km}$;
jet $3.628\times64.4 = \qty{233}{km}$.

\textbf{19.} $1.163\times29.4 - A\times\qty{233}{km} = 34.2 - 1.1 =
\qty{33.0}{\sverdrup}$; the excess, \qty{16}{\percent} of the interior
transport, returns south in the countercurrent just east of the jet.

\textbf{20.} $c = \sqrt{9.81\times4000} = \qty{198}{m/s}$; $L_R =
198/\num{8.37e-5} = \qty{2370}{km}$.

\textbf{21.} $198\times12.4206\times3600 = \qty{8860}{km}$.

\textbf{22.} Perimeter $2(6000 + 3336) = \qty{18670}{km}$, \qty{26.2}{h},
2.1 M$_2$ periods.

\textbf{23.} The circle is $2\pi a\cos\ang{35} = \qty{32790}{km}$, covered
in a lunar day, \qty{24.84}{h}: \qty{367}{m/s}, which needs $c^2/g =
\qty{13.7}{km}$ of water.

\textbf{24.} With a Rossby radius comparable to the width of the basin
and the M$_2$ wavelength comparable to its length, about one \omterm{def:b3:ocean-dynamics:amphidrome}{amphidromic
point}, near the centre of the basin; the high water turns anticlockwise
round it, keeping the coast on its right.

\textbf{25.} $\tfrac34\times0.357\times\cos^2\ang{35} = \qty{0.18}{m}$.

\textbf{26.} The western boundary current returns north the
\textbf{\qty[text-series-to-math=true]{29}{\sverdrup}} that the interior
carries south (\qty{29.4}{\sverdrup} by the calculation).
\end{solution}
